% Generated figure include blocks for Paper 7.

\begin{figure*}[t]
  \centering
  \includegraphics[width=0.98\textwidth]{../figures/fig1_clock_preserving.pdf}
  \caption{Area-preserving Hénon warping changes the potential geometry while preserving every configuration-sublevel area, and a fixed magnetic field changes the standard antiunitary symmetry while preserving every momentum-fiber area.  Both deformations therefore retain the same two-term Riemann--von Mangoldt mean clock.  The arithmetic prime-power gate remains open.}
  \label{fig:clock-preserving}
\end{figure*}

\begin{figure*}[t]
  \centering
  \includegraphics[width=0.96\textwidth]{../figures/fig2_equal_area_geometry.pdf}
  \caption{Configuration boundaries for two common energy levels.  The radial, centered (a=1.02), and high-distortion (a=6) domains have different geometry but exactly the same area at each energy by the determinant-one Hénon change of variables.}
  \label{fig:equal-area-geometry}
\end{figure*}

\begin{figure*}[t]
  \centering
  \includegraphics[width=0.98\textwidth]{../figures/fig3_classical_convergence.pdf}
  \caption{Classical diagnostics at (E=1000).  The radial finite-time exponent decays while its SALI remains order one; the centered nonlinear models retain order-one FTLE and roundoff-scale SALI.  An independent DOP853 implementation reproduces the (a=1.02) signal at (B=0) and retains it at (B=1), while the radial control has no joint flags.}
  \label{fig:classical-convergence}
\end{figure*}

\begin{figure*}[t]
  \centering
  \includegraphics[width=0.98\textwidth]{../figures/fig4_quantum_convergence.pdf}
  \caption{Failure-first quantum convergence audit.  The core (a=1.02) adjacent ratios stabilize across four meshes and (h^2) extrapolation, whereas the (a=6) branches remain sensitive at the available resolutions.  The first (0.04\to0.03) comparison failed the original one-percent level gate; later refinements are shown rather than substituted silently.}
  \label{fig:quantum-convergence}
\end{figure*}

\begin{figure}[t]
  \centering
  \includegraphics[width=0.72\textwidth]{../figures/fig5_ratio_cdfs.pdf}
  \caption{Empirical adjacent-ratio distributions over the frozen 138-ratio interior window.  The scalar spectrum is closer to the GOE surmise, while (B=1) moves toward the GUE surmise.  The sup distances are deterministic descriptive diagnostics, not i.i.d. goodness-of-fit tests.}
  \label{fig:ratio-cdfs}
\end{figure}

\begin{figure*}[t]
  \centering
  \includegraphics[width=0.94\textwidth]{../figures/fig6_magnetic_crossover.pdf}
  \caption{Magnetic response on the field grid frozen before execution; the uneven field values are placed at categorical horizontal positions for readability.  Every nonzero field lies above the scalar baseline, but the response is not monotone and no field is selected as an arithmetic optimum.  Independent coarse/fine pairs for the four new fields retain the direction of the response.}
  \label{fig:magnetic-crossover}
\end{figure*}

\begin{figure*}[t]
  \centering
  \includegraphics[width=0.94\textwidth]{../figures/fig7_smooth_clock.pdf}
  \caption{Exact-clock unfolding of the extrapolated (a=1.02) spectra.  The cumulative unfolded index follows unit mean spacing for both fields; the residual panel displays finite-window fluctuations only.  No Riemann ordinates enter this diagnostic.}
  \label{fig:smooth-clock}
\end{figure*}
